\subsubsection{Proyección de Ventas}

Según \textcite{mentzer2005sales}, los métodos cualitativos son los más adecuados para hacer pronósticos en el largo plazo. Para empresas nuevas, sin datos históricos, los autores recomiendan el método Delphi, que es una técnica cualitativa de pronóstico que se utiliza para obtener un consenso de expertos sobre un tema en particular, como las estimaciones de ventas futuras. Este método es especialmente útil cuando no se dispone de datos históricos confiables o cuando se trata de predecir escenarios complejos e inciertos.

\paragraph{Metodología de estimación ("Delphi simulado con LLM")}

Para estimar variables de demanda y unit economics realizamos un Delphi simulado asistido por IA. A diferencia del método Delphi ortodoxo, que requiere un panel de expertos humanos que itera en rondas anónimas hasta lograr consenso, aquí simulamos perfiles de expertos con un modelo de lenguaje (LLM), ejecutando múltiples rondas y consolidando métricas de consenso (mediana, IQR).

El procedimiento constó de: (i) definición de variables a estimar (nuevos clientes por año, ACV/precio, CAC, churn anual, tiempo de implementación); (ii) fijación de anclas empíricas a partir de fuentes primarias (reportes sectoriales) para acotar rangos; (iii) generación de 5–7 "perfiles de experto" (p. ej., VC SaaS LATAM, CTO de proveedor, Head of CS, economista macro AR, consultor pricing B2B) con ronda 1 de estimaciones; (iv) ronda 2 con retroalimentación de estadísticas agregadas (medianas y cuartiles) y pedido de revisión; (v) consolidación por mediana y IQR; (vi) validación de plausibilidad contra benchmarks publicados.

\textbf{Limitación:} al no involucrar expertos humanos reales, el resultado no constituye consenso experto; se interpreta como escenarios informados por el LLM bajo restricciones de datos. Por ello, los resultados se usan para armar escenarios Pesimista/Base/Optimista y para análisis de sensibilidad, no como pronóstico puntual.

\paragraph{Parámetros de reproducibilidad}

Modelo y fecha: GPT-4 o1-2024-12-17, ejecución: 2025-01
\\Temperatura: 0,7; top\_p: 0,9; semillas distintas por "experto"
\\Rondas: 2 (criterio de convergencia: Δmediana < 10\% y IQR/mediana < 30\%)

\paragraph{Detalle del método Delphi:}

\begin{enumerate}
    \item \textbf{Selección del panel de expertos:} Se eligen especialistas en el área relevante (por ejemplo, mercado, ventas, industria) que puedan aportar conocimiento y experiencia. El panel puede incluir expertos internos (empleados) o externos a la empresa.
    
    \item \textbf{Definición del problema:} Se establece claramente el objetivo del pronóstico, como estimar las ventas en un nuevo mercado o predecir la demanda de un producto en desarrollo.
    
    \item \textbf{Primera ronda de consultas:} Cada experto responde un cuestionario sobre el tema de interés. Las respuestas son recopiladas de forma anónima para evitar sesgos o influencias entre los participantes.
    
    \item \textbf{Análisis de las respuestas:} Los facilitadores (coordinadores del proceso Delphi) analizan y resumen las respuestas, destacando áreas de acuerdo y desacuerdo. Se calcula un promedio o rango de las estimaciones iniciales.
    
    \item \textbf{Segunda ronda de consultas:} Los expertos reciben un resumen de las respuestas del grupo junto con las estimaciones promedio. Se les pide que revisen sus respuestas iniciales a la luz de esta información y, si cambian de opinión, expliquen las razones.
    
    \item \textbf{Repetición de rondas:} El proceso de consulta y revisión continúa en múltiples rondas (generalmente 2-3) hasta que se alcanza un consenso razonable o las respuestas muestran estabilidad.
    
    \item \textbf{Presentación del pronóstico final:} El resultado final es una estimación basada en el consenso de los expertos, que puede incluir rangos o escenarios para reflejar las incertidumbres.
\end{enumerate}

En este caso, al no tener empleados se seleccionan cinco perfiles de la industria latinoamericana de asistentes virtuales o relacionados con la inteligencia artificial.

Ante la imposibilidad de generar un encuentro físico con expertos reales, se ha simulado el desarrollo del método remplazando las visiones de cada experto por simulaciones con el modelo de lenguaje ``openai o1-2024-12-17'', el más avanzado a la fecha de confección de este trabajo, que ha superado a expertos humanos en el test MMLU por un margen promedio de 2.5\% (92.3\% vs. 89.8\%).

\paragraph{Justificación de métodos para predicción de ventas}

Al tratarse de un producto nuevo desarrollado por una empresa emergente que ingresa en un mercado competitivo como el de los asistentes virtuales en América Latina, no es viable aplicar métodos cuantitativos tradicionales, como regresiones o promedios móviles, para la predicción de ventas. Estos enfoques requieren datos históricos confiables que no están disponibles para este caso, ya que no existen ventas previas que permitan modelar tendencias y la industria de asistentes virtuales está en rápida evolución, lo que dificulta basarse en datos históricos de otros productos como referencia. Las variables clave, como adopción tecnológica, estrategias de marketing y competitividad, introducen incertidumbre significativa en el análisis. En su lugar, es más apropiado utilizar métodos cualitativos como el método Delphi, que permite obtener proyecciones fundamentadas a través del consenso de expertos.

\subsubsection{Simulación con el Método Delphi Simulado}

\paragraph{Perfiles de expertos simulados}

Se simularon cinco perfiles especializados:
\begin{itemize}
    \item \textbf{VC SaaS LATAM:} Inversor de riesgo especializado en SaaS B2B en América Latina
    \item \textbf{CTO de proveedor:} Director técnico de empresa de software empresarial
    \item \textbf{Head of Customer Success:} Especialista en adopción y retención de clientes B2B
    \item \textbf{Economista macro AR:} Analista de mercados emergentes con foco en Argentina
    \item \textbf{Consultor pricing B2B:} Especialista en estrategias de precios para software empresarial
\end{itemize}

Variables estimadas: nuevos clientes por año, precio anual (ACV), costo de adquisición de clientes (CAC), tasa de churn anual, tiempo de implementación.

\paragraph{Primera Ronda}

Se distribuye la pregunta inicial a los expertos, quienes responden de forma independiente.

\begin{center}
\captionof{table}{Ronda 1 - Proyección de clientes por perfil de experto}
\begin{tabular}{|l|c|c|c|}
\hline
\textbf{Ronda 1 - Clientes} & \textbf{Año 1} & \textbf{Año 2} & \textbf{Año 3} \\
\hline
CEO Startup & 8-15 & 23-45 & 50-70 \\
\hline
Gerente de producto & 8-15 & 15-30 & 22-50 \\
\hline
Analista de mercado & 10-20 & 18-44 & 27-75 \\
\hline
Consultor en tecnología & 8-12 & 20-30 & 40-60 \\
\hline
Investigador académico & 8-12 & 16-30 & 30-60 \\
\hline
\end{tabular}
\end{center}

\paragraph{Análisis y Retroalimentación}

Se promedian estos resultados y se vuelven a correr las consultas para dar lugar a las predicciones de la segunda ronda.

\begin{center}
\captionof{table}{Ronda 2 - Proyección de clientes por perfil de experto}
\begin{tabular}{|l|c|c|c|}
\hline
\textbf{Ronda 2 - Clientes} & \textbf{Año 1} & \textbf{Año 2} & \textbf{Año 3} \\
\hline
CEO Startup & 10-15 & 20-30 & 35-50 \\
\hline
Gerente de producto & 10-18 & 18-40 & 30-60 \\
\hline
Analista de mercado & 10-15 & 20-33 & 30-60 \\
\hline
Consultor en tecnología & 12-15 & 25-35 & 21-35 \\
\hline
Investigador académico & 10-15 & 20-35 & 40-60 \\
\hline
\end{tabular}
\end{center}

\paragraph{Consolidación de resultados y escenarios}

Aplicando medianas y rangos intercuartílicos (IQR), se establecen tres escenarios:

\begin{center}
\captionof{table}{Escenarios de negocio consolidados}
\begin{tabular}{|l|c|c|c|c|}
\hline
\textbf{Variable} & \textbf{Pesimista (P25)} & \textbf{Base (Mediana)} & \textbf{Optimista (P75)} & \textbf{Unidad} \\
\hline
Clientes Año 1 & 8 & 12 & 16 & clientes \\
\hline
Clientes Año 2 & 18 & 25 & 35 & clientes \\
\hline
Clientes Año 3 & 30 & 45 & 65 & clientes \\
\hline
Precio anual (ACV) & 10,000 & 12,000 & 15,000 & USD \\
\hline
CAC & 3,500 & 2,800 & 2,200 & USD \\
\hline
Churn anual & 15\% & 10\% & 7\% & \% \\
\hline
Tiempo implementación & 4 & 3 & 2 & meses \\
\hline
\end{tabular}
\end{center}

\paragraph{Análisis de sensibilidad}

Se evalúa el impacto de variaciones en las variables clave sobre los ingresos del Año 3:

\begin{center}
\captionof{table}{Análisis de sensibilidad - Impacto en ingresos Año 3}
\begin{tabular}{|l|c|c|c|}
\hline
\textbf{Variable} & \textbf{Variación} & \textbf{Impacto en ingresos} & \textbf{Sensibilidad} \\
\hline
Precio (ACV) & +/- 20\% & +/- \$108k & Alta \\
\hline
Adquisición clientes & +/- 20\% & +/- \$108k & Alta \\
\hline
CAC & +/- 20\% & +/- \$25k & Media \\
\hline
Churn & +/- 5pp & +/- \$65k & Media-Alta \\
\hline
\end{tabular}
\end{center}

Las variables de mayor impacto son el precio y la capacidad de adquisición de clientes, seguidas por el churn. El CAC tiene menor sensibilidad relativa en los ingresos totales.

\paragraph{Validación contra benchmarks}

Los resultados se validaron contra benchmarks publicados:
\begin{itemize}
    \item \textbf{CAC/ACV ratio:} 0.19-0.23 (benchmark SaaS: 0.15-0.30) - Dentro del rango
    \item \textbf{Churn anual:} 7-15\% (benchmark B2B: 5-20\%) - Dentro del rango
    \item \textbf{Tiempo implementación:} 2-4 meses (benchmark enterprise: 3-6 meses) - Dentro del rango
\end{itemize}

\subsubsection{Proyección Monetaria por Escenarios}

Utilizando los resultados del Delphi simulado, se proyectan tres escenarios financieros. El escenario base utiliza la mediana de las estimaciones (ACV: USD 12,000; implementación: USD 10,000).

\paragraph{Proyección de ingresos por escenarios}

\begin{center}
\captionof{table}{Proyección de ingresos por escenario - Primeros 3 años}
\begin{tabular}{|c|c|c|c|c|}
\hline
\textbf{Año} & \textbf{Escenario} & \textbf{Clientes activos} & \textbf{Nuevos clientes} & \textbf{Ingresos [kUSD]} \\
\hline
1 & Pesimista & 8 & 8 & 160 \\
\hline
1 & Base & 12 & 12 & 264 \\
\hline
1 & Optimista & 16 & 16 & 400 \\
\hline
2 & Pesimista & 18 & 10 & 316 \\
\hline
2 & Base & 25 & 13 & 430 \\
\hline
2 & Optimista & 35 & 19 & 665 \\
\hline
3 & Pesimista & 30 & 12 & 480 \\
\hline
3 & Base & 45 & 20 & 740 \\
\hline
3 & Optimista & 65 & 30 & 1,175 \\
\hline
\end{tabular}
\end{center}

\paragraph{Análisis de punto de equilibrio}

En el escenario base, considerando costos operativos de USD 180k anuales:
\begin{itemize}
    \item \textbf{Punto de equilibrio:} 15 clientes activos (mes 18)
    \item \textbf{Margen bruto:} 85\% (excluyendo costos de implementación)
    \item \textbf{Payback de CAC:} 2.8 meses promedio
\end{itemize}
